% =====================================================================
%  Relatório: Atividade 05 (Maix que um Circuito, jogo 3D)
%  MCCC007-23 Computação Gráfica (UFABC), Prof. João Paulo Gois
%
%  Figuras: python capturas.py   (abre o jogo e gera figuras/*.jpg)
%  Compilar duas vezes:  pdflatex relatorio.tex && pdflatex relatorio.tex
%
%  \preencher{...} aparece em VERMELHO no PDF: algo que os autores ainda
%  precisam confirmar. Antes de entregar: grep -n 'preencher{' relatorio.tex
% =====================================================================
\documentclass[a4paper,11pt]{article}

\usepackage[T1]{fontenc}
\usepackage[utf8]{inputenc}
\usepackage[brazilian]{babel}
\usepackage{lmodern}
\usepackage{microtype}
\usepackage[margin=2.3cm]{geometry}
\usepackage{enumitem}
\usepackage{amsmath,amssymb}
\usepackage{graphicx}
\graphicspath{{figuras/}}
\usepackage{float}
\usepackage{caption}
\usepackage{subcaption}
\usepackage{booktabs}
\usepackage{tabularx}
\newcolumntype{Y}{>{\raggedright\arraybackslash}X}
\usepackage{xcolor}
\usepackage[hidelinks]{hyperref}

\newcommand{\preencher}[1]{\textcolor{red}{\textbf{[}#1\textbf{]}}}
\newcommand{\tecla}[1]{\fbox{\footnotesize\ttfamily #1}}
\newcommand{\arq}[1]{\texttt{#1}}

\setlist{itemsep=2pt,topsep=4pt}

\title{\textbf{Maix que um Circuito}\\[2pt]
  \large Jogo 3D em ModernGL sobre o problema do corte máximo\\[2pt]
  \normalsize MCCC007-23 Computação Gráfica, Atividade 05, Prof.\ João Paulo Gois}
\author{Lucas Ataide (RA 11202421363) \and Gabriel Frigo (RA 11202321002)}
\date{24 de novembro de 2026\\[2pt]
  \small Repositório: \url{https://github.com/lucas-ata1de/maix-que-um-circuito}}

\begin{document}
\maketitle

% =====================================================================
\section{O jogo}
% =====================================================================
A arena tem $n=16$ torres de transmissão, que são os vértices de um grafo
aleatório; os fios entre elas são as arestas. Cada torre pertence a um de dois
lados (azul ou vermelho). Uma aresta está \emph{no corte} quando liga torres de
lados diferentes: nesse caso o fio faísca em amarelo; caso contrário, fica cinza
e quase parado. O jogador anda pela arena e clica nas torres para trocá-las de
lado. Vence quem levar o número de arestas no corte até a \emph{meta} antes de o
cronômetro de 60\,s zerar; se o tempo acabar, perde.

O tema vem da disciplina de Algoritmos Probabilísticos: a partição inicial é
sorteada pelo algoritmo aleatorizado de $\tfrac12$-aproximação para o corte
máximo (sorteia cada vértice com probabilidade $\tfrac12$ até obter um corte com
pelo menos $m/2$ arestas), e a meta da fase é o melhor corte encontrado por
busca local gulosa com 10 reinícios aleatórios. Como a busca local só faz inversões
de um vértice, o mesmo movimento do jogador, a meta é sempre alcançável. A fase
nunca começa já vencida: o sorteio é refeito enquanto o corte inicial for maior
ou igual à meta.

\paragraph{Controles.} \tecla{W}\tecla{A}\tecla{S}\tecla{D} andam,
\tecla{Espaço}/\tecla{Shift} sobem e descem (câmera livre);
\tecla{T} alterna a câmera; \tecla{Enter} inicia; \tecla{P} pausa;
\tecla{R} reinicia com um novo grafo; \tecla{1}/\tecla{2} ligam e desligam a luz direcional e a pontual móvel; \tecla{3} liga e desliga todas as luzes dos choques, \tecla{4} só as das torres e \tecla{5} só as das arestas; \tecla{Y} liga e desliga o VSync; \tecla{C}/\tecla{V}
desligam e ligam o descarte de faces; \tecla{Esc} fecha. Mouse: clique esquerdo
numa torre troca o lado dela; arrastar com o botão direito gira a câmera; a roda
aproxima e afasta. A mesma lista está no início de \arq{main.py}.

% =====================================================================
\section{Arquitetura}
% =====================================================================
\begin{table}[H]
\centering\small
\begin{tabularx}{\linewidth}{@{}lYl@{}}
\toprule
\textbf{Módulo} & \textbf{Responsabilidade} & \textbf{Origem} \\
\midrule
\arq{main.py} & Janela GLFW, callbacks, laço de jogo ($\Delta t$, FPS, VSync), matrizes de modelo de cada objeto e ordem das passadas de desenho. & novo \\
\arq{jogo.py} & Classe \texttt{EstadoJogo}: estado da partida, cronômetro, sorteio das posições, \emph{picking} (raio $\times$ caixa) e colisão esfera $\times$ caixa. & novo \\
\arq{grafo.py} & Classe \texttt{Graph}, geração $G(n,p)$ conexa, corte aleatorizado $\ge m/2$ e busca local (meta). & novo \\
\arq{camera.py} & Classe \texttt{Camera}: modos livre e aéreo, matrizes de visão (LookAt) e projeção. & adaptado da apostila \\
\arq{renderizador.py} & Classe \texttt{Renderizador}: única que conhece shaders, VAOs e uniforms; três passadas (objetos, linhas, placar). & novo \\
\arq{material.py} & \texttt{Material}, \texttt{LuzDirecional}, \texttt{LuzPontual}, \texttt{CenaLuzes}, \texttt{ConfigLuzesChoque} (luzes dos choques) e conversão sRGB$\to$linear. & adaptado da apostila \\
\arq{malha.py}, \arq{texturas.py} & Leitura do OFF, normais, mapeamento cilíndrico com costura; carga de imagens em \texttt{GL\_SRGB8\_ALPHA8} com mipmaps. & visualizador da apostila \\
\arq{geometria.py} & Geometrias procedurais (piso, cubo, anel, segmento, quadrado do placar) e enquadramento da torre por matriz. & novo \\
\bottomrule
\end{tabularx}
\caption{Módulos do projeto.}
\end{table}

\paragraph{Como as classes se comunicam.} \texttt{main} cria uma instância de
cada classe e as conecta pelo laço. Os callbacks de teclado e mouse traduzem
eventos em chamadas: o clique chama \texttt{Camera.raio} para obter o raio do
cursor e \texttt{EstadoJogo.torre\_no\_raio} para achar a torre atingida, e então
\texttt{EstadoJogo.inverter\_torre}, que usa o \texttt{Graph} para trocar o lado
e recontar o corte. A cada quadro, \texttt{main} chama \texttt{Camera.mover},
\texttt{EstadoJogo.empurrar\_esfera} (colisão) e \texttt{EstadoJogo.atualizar(dt)}
e depois entrega ao \texttt{Renderizador} as matrizes e os \texttt{Material}
de cada objeto. A lógica do jogo (\arq{jogo.py}, \arq{grafo.py}) não importa
\texttt{moderngl}: só o \texttt{Renderizador} fala com a GPU.

\paragraph{Shaders.} São dois programas.
\arq{cena.vert}/\arq{cena.frag} desenham os objetos iluminados: o vertex shader
leva posição e normal ao espaço da câmera; o fragment shader avalia Blinn--Phong para a luz direcional, para a luz pontual móvel e para as luzes dos choques (Seção~5) e multiplica $\kappa_a$ e $\kappa_d$ pela cor da textura, cujas coordenadas o vertex shader multiplica por \texttt{u\_uv\_escala}.
\arq{linha.vert}/\arq{linha.frag} desenham as linhas e o placar com cor sólida.
Nas arestas, o vertex shader desloca os vértices internos de um segmento
unitário com senos que dependem do tempo, multiplicados pelo envelope
$\sin(\pi z)$, que vale zero nas pontas e mantém as pontas presas no topo das
torres; isso produz a faísca sem reescrever o VBO.

\paragraph{Reaproveitamento.} Do jogo 2D (Atividade 03) vem o padrão da
janela e do laço \emph{hints} de contexto 4.0 core, medição de $\Delta t$ com
passo máximo, FPS no título e VSync por tecla. Do visualizador da apostila vêm
\arq{malha.py} e \arq{texturas.py} (usados como estão) e a base de
\arq{camera.py} e \arq{material.py}. Na câmera foram criados os dois modos, a
navegação por yaw/pitch e o método \texttt{raio}; no material, os presets do
jogo. O resto é do grupo. O tema do jogo é novo, não uma extensão do jogo 2D.

\paragraph{VBOs.} Todos os VBOs (torre com IBO, piso, cubo, anel, segmento e
quadrado do placar) e as texturas são criados uma única vez, antes do laço.
Cada torre, borda, anel, aresta e barra é o mesmo VBO desenhado com outra matriz
de modelo. Nenhum VBO é reescrito durante o jogo, então não há exceção a
justificar. Primitivas usadas: \texttt{TRIANGLES} (objetos),
\texttt{LINE\_STRIP} (arestas), \texttt{LINE\_LOOP} (anéis e molduras do
placar) e \texttt{TRIANGLE\_FAN} (barras do placar).

% =====================================================================
\section{Laço de jogo e animação}
% =====================================================================
A cada quadro, $\Delta t$ é a diferença entre dois \texttt{glfw.get\_time()}
consecutivos, limitada a \texttt{DT\_MAX} $=0{,}05$\,s. Toda velocidade está em
unidades por segundo e é multiplicada por $\Delta t$: a câmera anda a 6\,u/s, o
cronômetro desconta $\Delta t$, e os anéis giram $2{,}2$\,rad/s. Por isso a
velocidade do jogo não depende da taxa de quadros: com 60 ou com 280 quadros
por segundo, a soma dos $\Delta t$ em um segundo é a mesma, e o deslocamento
também. A tecla \tecla{Y} chama
\texttt{glfw.swap\_interval(0 ou 1)} e o FPS aparece no título, recalculado a
cada 0,25\,s.

\begin{table}[H]
\centering\small
\begin{tabular}{@{}lc@{}}
\toprule
\textbf{Modo} & \textbf{FPS observado} \\
\midrule
VSync ligado & 60,1 \\
VSync desligado & 262,2 (entre 260 e 370 em outras execuções) \\
\bottomrule
\end{tabular}
\caption{Taxas medidas pelo próprio contador do jogo, lidas do título pelo
script de capturas (janela 900$\times$900, placa integrada).}
\end{table}

\paragraph{Animação autônoma.} Três elementos se movem sozinhos, em função do
tempo de animação: os anéis no topo das torres giram em torno de $y$, as arestas
faíscam a luz pontual orbita a arena num círculo de raio 8 com uma leve oscilação vertical, e as luzes dos choques cintilam (a intensidade oscila com senos do tempo).

\paragraph{Estados e pausa.} Os estados são \texttt{START}, \texttt{PLAYING},
\texttt{PAUSED}, \texttt{VICTORY} e \texttt{GAME\_OVER}. \tecla{Enter} inicia,
\tecla{P} pausa e retoma, e \tecla{R} sorteia uma nova fase sem fechar o
programa. Na pausa, nada se move: \texttt{EstadoJogo.atualizar} só desconta tempo
em \texttt{PLAYING}, a câmera só anda em \texttt{START}/\texttt{PLAYING} e o
tempo de animação (\texttt{tempo\_anim}), que comanda anéis, faíscas e luzes, só avança fora da pausa. No fim de jogo a cena continua animada, mas o cronômetro, os cliques e a câmera param. O tempo pausado não entra no $\Delta t$:
o instante anterior é atualizado em todo quadro, inclusive durante a pausa,
então o primeiro $\Delta t$ depois da retomada é de um quadro só, e ainda é
limitado por \texttt{DT\_MAX}. No fim, o título mostra
``GAME OVER, tente novamente aperte R'' ou
``Ganhou! Tente novamente apertando R. Sobrou $x$ segundos.''

% =====================================================================
\section{Transformações e câmera}
% =====================================================================
\paragraph{Matrizes de modelo.} Cada objeto recebe sua própria matriz
$M$, sempre na ordem escala, depois rotação, depois translação (multiplicação
da direita para a esquerda):
\begin{itemize}
\item \textbf{Torre} (arquivo OFF): $M = T(x,0,z)\,M_{\text{enq}}$, com
  $M_{\text{enq}} = S(4/h)\,T(-c_x,-y_{\min},-c_z)$, onde $h$ é a altura da caixa
  envolvente e $c$ o seu centro. Os vértices do arquivo vão para o VBO sem
  alteração; o enquadramento (pés em $y=0$, altura 4) é feito só pela matriz.
\item \textbf{Bordas}: $T\,S$ (frente e fundo) e $T\,R_y(90^\circ)\,S$ (laterais) sobre o cubo unitário, com escala \emph{não uniforme} $S=(24{,}4;\,0{,}4;\,0{,}4)$ ou $(23{,}6;\,0{,}4;\,0{,}4)$. A rotação faz a textura correr ao longo do comprimento nas quatro bordas.
\item \textbf{Anéis}: $T(\text{topo})\,R_y(2{,}2\,t + 0{,}7\,u)$.
\item \textbf{Arestas}: $T(A)\,R_y(\operatorname{atan2}(d_x,d_z))\,S(1,1,\lVert B-A\rVert)$,
  que leva o segmento $(0,0,0)\to(0,0,1)$ até o segmento $A\to B$.
\item \textbf{Placar}: $T(x,y,0)\,S(\ell,a,1)$ sobre o quadrado $[0,1]^2$.
\end{itemize}

\paragraph{Visão.} As duas câmeras usam \texttt{glm.lookAt}.
No modo \textbf{livre}, o olho é \texttt{pos} e o alvo é
\texttt{pos + frente}, com
$\text{frente} = (\cos\phi\cos\theta,\ \sin\phi,\ \cos\phi\sin\theta)$, sendo
$\theta$ o yaw e $\phi$ o pitch (limitado a $\pm85^\circ$ para a direção de
visão nunca ficar paralela ao vetor \emph{up}). WASD movem no plano horizontal,
o botão direito altera yaw e pitch e a roda avança. No modo \textbf{aéreo}, o
olho fica em $(0,\,h,\,0{,}2)$ olhando para a origem, com $h$ ajustado pela roda
entre 4 e 60; o pequeno deslocamento em $z$ evita a direção de visão paralela
ao vetor \emph{up}. \tecla{T} alterna os modos.

\paragraph{Projeção.} Perspectiva, com \texttt{glm.perspective} de
$45^\circ$. O \textbf{near} é $0{,}1$: menor que o raio de colisão da câmera
($0{,}5$), de modo que o plano near nunca entra no chão nem nas torres. O
\textbf{far} é 150: a arena tem $24\times24$ e, com a câmera aérea no alto
(60), o canto mais distante da arena fica a cerca de 62 unidades; 150 deixa
folga também para a câmera livre vista de longe, sem desperdiçar a precisão do
depth buffer (a razão far/near é 1500).

\paragraph{Redimensionamento.} \texttt{glfw.set\_framebuffer\_size\_callback}
ajusta \texttt{ctx.viewport} para o novo tamanho do framebuffer e atualiza a
razão de aspecto da câmera ($\text{largura}/\text{altura}$), usada pela
projeção no quadro seguinte. A janela é redimensionável.

\paragraph{Matriz normal.} O \texttt{Renderizador} calcula, na CPU, para cada
objeto, $N = \big((V M)_{3\times3}^{-1}\big)^{T}$ e envia como
\texttt{u\_normal\_matrix}. Ela é necessária por causa das bordas, que têm escala
não uniforme: aplicar $VM$ diretamente à normal a incliná-la e a deixaria fora da
perpendicular à superfície.

% =====================================================================
\section{Iluminação e materiais}
% =====================================================================
O modelo é \textbf{Blinn--Phong} avaliado no fragment shader (sombreamento de
Phong). Para cada luz ativa:
\[
I = c_{\text{luz}}\big(f_a\,\kappa_a\,t
  + \kappa_d\,t\,\max(N\!\cdot\!L,0)
  + \kappa_s\,\max(N\!\cdot\!H,0)^{\alpha}\big),
\qquad H = \frac{L+V}{\lVert L+V\rVert},
\]
em que $t$ é a cor da textura e $f_a$ a fração ambiente da luz. A luz pontual
ainda é multiplicada pela atenuação $1/(k_c + k_l d + k_q d^2)$. O especular
só é somado quando $N\cdot L>0$.

\paragraph{Espaço.} Tudo é calculado no \textbf{espaço da câmera}: o vertex
shader envia posição e normal multiplicadas por $V M$; na CPU, a direção da luz
direcional é multiplicada por $V_{3\times3}$ (vetor, sem translação) e a posição
da luz pontual por $V$ com $w=1$ (ponto). Assim o vetor de visão é
simplesmente $V = -p$.

\begin{table}[H]
\centering\small
\begin{tabularx}{\linewidth}{@{}lY@{}}
\toprule
\textbf{Fonte} & \textbf{Parâmetros} \\
\midrule
Direcional (\tecla{1}) & direção $(-0{,}4,-1,-0{,}35)$, cor $(1{,}20;1{,}22;1{,}26)$, $f_a=0{,}28$. \\
Pontual móvel (\tecla{2}) & órbita de raio 8 a $0{,}8$\,rad/s, altura $4{,}5+0{,}6\sin(1{,}6t)$; cor $(4{,}40;3{,}45;2{,}15)$, $f_a=0{,}12$; $k_c=1$, $k_l=0{,}045$, $k_q=0{,}010$. \\
Choques das torres (\tecla{3}/\tecla{4}) & uma fonte pontual por torre, 0,25 acima do topo; cor azul $(0{,}40;1{,}00;2{,}25)$ ou vermelha $(1{,}75;0{,}36;0{,}38)$ conforme o lado da torre, vezes $0{,}9+0{,}1\sin(7{,}5t+1{,}9u)$; $k_c=1$, $k_l=0{,}30$, $k_q=0{,}20$, $f_a=0{,}08$, alcance $R=4{,}8$. \\
Choques das arestas (\tecla{3}/\tecla{5}) & fonte \emph{linear}: a luz parte do ponto do segmento $AB$ mais próximo do fragmento; cor amarela $(1{,}25;1{,}02;0{,}18)$ no corte e azul-escura $(0{,}12;0{,}16;0{,}24)$ fora, com intensidade $(1-a)+a\sin(\ldots)$, $a=0{,}15$ ou $0{,}06$; $k_c=1$, $k_l=0{,}48$, $k_q=0{,}32$, $f_a=0{,}03$, alcance $R=3{,}6$. \\
\bottomrule
\end{tabularx}
\caption{Fontes de luz. Todas ficam ativas ao mesmo tempo e podem ser alternadas por tecla. O shader aceita até 24 choques de torre e 48 de aresta; se houver mais arestas, as do corte têm prioridade.}
\end{table}

\paragraph{Luzes dos choques.} Com 16 torres e dezenas de arestas, somar
todas as fontes com a atenuação comum clarearia a arena inteira. Por isso os
choques usam uma atenuação com \emph{alcance finito} $R$:
\[
a(d) = \frac{\big(1-(d/R)^2\big)^2}{k_c + k_l d + k_q d^2}\ \text{ se } d<R,
\qquad a(d)=0 \text{ caso contrário},
\]
e o shader pula a fonte quando $a=0$. Na luz de aresta, $d$ é a distância até
o ponto $A + t(B-A)$ com
$t=\operatorname{clamp}\!\big((p-A)\cdot(B-A)/\lVert B-A\rVert^2,\,0,\,1\big)$,
o mesmo ponto mais próximo usado na colisão. As posições $A$, $B$ e as das
torres são levadas ao espaço da câmera na CPU, como a luz pontual.

\begin{table}[H]
\centering\small
\begin{tabular}{@{}lllll@{}}
\toprule
\textbf{Material} & $\kappa_a$ & $\kappa_d$ & $\kappa_s$ & $\alpha$ \\
\midrule
Piso (textura) & $(1{,}80;1{,}80;1{,}80)$ & $(3{,}60;3{,}60;3{,}60)$ & $(0{,}12;0{,}12;0{,}14)$ & 32 \\
Torre azul (textura) & $(0{,}32;0{,}65;1{,}25)$ & $(0{,}55;1{,}25;2{,}35)$ & $(0{,}50;0{,}60;0{,}70)$ & 64 \\
Torre vermelha (textura) & $(0{,}68;0{,}24;0{,}24)$ & $(1{,}40;0{,}44;0{,}42)$ & $(0{,}60;0{,}50;0{,}46)$ & 64 \\
Bordas (textura) & $(0{,}16;0{,}16;0{,}16)$ & $(0{,}46;0{,}46;0{,}46)$ & $(0{,}04;0{,}04;0{,}04)$ & 12 \\
\bottomrule
\end{tabular}
\caption{Materiais (coeficientes em espaço linear). O lado de cada torre é mostrado pela troca de material. Coeficientes acima de 1 compensam texturas escuras.}
\end{table}

Dos itens de iluminação, o jogo cumpre os cinco: três tipos de fonte (direcional, pontual com atenuação e linear), uma fonte
que se move, quatro materiais, matriz normal com objeto de escala não uniforme e
todos os vetores no mesmo espaço (o da câmera).

% =====================================================================
\section{Texturas}
% =====================================================================
\begin{table}[H]
\centering\small
\begin{tabularx}{\linewidth}{@{}llY@{}}
\toprule
\textbf{Imagem} & \textbf{Objeto} & \textbf{Fonte e licença} \\
\midrule
\arq{rubber\_tiles\_diff\_1k.jpg} & piso & Poly Haven, ``Rubber Tiles'', \url{https://polyhaven.com/a/rubber_tiles}, CC0 1.0 Universal (Domínio Público) \\
\arq{rusty\_metal\_04\_diff\_1k.jpg} & torres & Poly Haven, ``Rusty Metal 04'', \url{https://polyhaven.com/a/rusty_metal_04}, CC0 1.0 Universal (Domínio Público) \\
\arq{white\_stucco\_02\_diff\_1k.jpg} & bordas & Poly Haven, ``White Stucco 02'', \url{https://polyhaven.com/a/white_stucco_02}, CC0 1.0 Universal (Domínio Público) \\
\bottomrule
\end{tabularx}
\caption{Texturas de cor (mapas difusos de $1024\times1024$). Todas as texturas do Poly Haven são distribuídas sob licença CC0 1.0 Universal (\url{https://polyhaven.com/license}).}
\end{table}

\paragraph{Coordenadas de textura.} No \textbf{piso}, são geradas junto com a
geometria, de $(0,0)$ a $(12,12)$, de modo que a imagem se repete $12\times12$
vezes. Nas \textbf{torres}, são calculadas por \textbf{mapeamento cilíndrico}
em torno de $y$ ($u$ pelo ângulo em torno do eixo, $v$ pela altura). A
\textbf{costura} é tratada: num triângulo que cruza a linha $u=0/u=1$ (variação
de $u$ maior que meia volta), os vértices com $u<0{,}5$ ganham uma cópia com
$u+1$ e o triângulo passa a usar a cópia. A malha passa de 1247 para 1288
vértices, e some a faixa em que a textura inteira apareceria espremida num único triângulo. Nas \textbf{bordas}, as coordenadas de 0 a 1 de cada face do cubo são multiplicadas no vertex shader por \texttt{u\_uv\_escala} $=(20,1)$, e a textura se repete 20 vezes ao longo do comprimento.

\paragraph{Empacotamento e filtros.} As três texturas usam \texttt{GL\_REPEAT}, necessário no piso e nas bordas, onde as coordenadas passam de 1, e na torre, por causa da costura com $u+1$. Os mipmaps são gerados na inicialização
(\texttt{build\_mipmaps}) e a minificação é trilinear
(\texttt{LINEAR\_MIPMAP\_LINEAR}, com filtragem anisotrópica de até 8 amostras);
a magnificação é linear. Sem os mipmaps, o padrão do piso tremeria ao longe.

\paragraph{sRGB.} Na \textbf{leitura}, as imagens são carregadas em
\texttt{GL\_SRGB8\_ALPHA8}: a GPU converte cada texel de sRGB para linear antes do
shader. A iluminação é somada em espaço linear, e as cores das linhas e do
placar também são convertidas para linear (na CPU, uma vez, com
\texttt{srgb\_para\_linear}). Na \textbf{escrita}, a cor é codificada em sRGB uma
única vez, pelo estado \texttt{GL\_FRAMEBUFFER\_SRGB}, com uma janela pedida com
\texttt{SRGB\_CAPABLE}. A cor de fundo já está em sRGB; por isso o estado é
desligado durante o \texttt{clear} e religado logo depois, evitando que ela seja
codificada duas vezes.

\paragraph{Modelo de terceiros.} A torre vem do modelo ``High Voltage Transmission Line Tower (tileable)'', de \emph{mednios} no Sketchfab (\url{https://skfb.ly/ot8Sx}), sob a licença Creative Commons Atribuição 4.0 (\url{http://creativecommons.org/licenses/by/4.0/}). O arquivo OBJ foi convertido para OFF fora do programa por
um script Python. Os seis cabos e os isoladores, que se estendiam por 44
unidades, foram descartados, e só a estrutura da torre foi mantida
(1247 vértices, 2031 triângulos). As coordenadas restantes não foram alteradas.

% =====================================================================
\section{Colisões}
% =====================================================================
\paragraph{Formas.} A câmera livre é o corpo do jogador e é tratada como uma
\textbf{esfera} de raio $r=0{,}5$. Cada torre é envolvida por uma \textbf{caixa
alinhada aos eixos} (AABB) com meia largura 1 e altura 4,1, e o chão e as
paredes da arena são planos.

\paragraph{Teste e reação (esfera $\times$ AABB).} O ponto da caixa mais
próximo do centro $p$ é $q = \operatorname{clamp}(p, b_{\min}, b_{\max})$,
coordenada a coordenada. Há colisão se $0 < \lVert p-q\rVert < r$. A reação
empurra a esfera para fora ao longo da normal:
$p \leftarrow q + (p-q)\,r/\lVert p-q\rVert$. Só a parte do movimento que entra
na torre é desfeita; a parte tangente continua, então a câmera desliza pela
lateral da torre em vez de travar. Depois disso, um \emph{clamp} mantém
$p_y \ge r$ (chão) e $|p_x|,|p_z| \le 12-r$ (paredes). Como as caixas têm
altura, a colisão é de fato 3D: a câmera pode passar por cima das torres.

\paragraph{Tunelamento.} O teste só funciona enquanto o centro da esfera não
entra na caixa (quando entra, $p=q$ e não há direção para empurrar). Isso é
garantido por uma regra: \emph{todo deslocamento da câmera é menor que $r$ e a
colisão é resolvida logo depois de cada deslocamento}. Andando, o maior passo
é $6\,\text{u/s} \times 0{,}05\,\text{s} = 0{,}3 < 0{,}5$, graças ao passo
máximo de $\Delta t$; sem ele, um travamento de 1\,s faria a câmera saltar 6
unidades e atravessar a torre. Com a roda, cada evento move $0{,}4 < 0{,}5$ e a
colisão é resolvida dentro do próprio callback, porque vários eventos podem
chegar num mesmo \texttt{poll\_events}.

\paragraph{Seleção pelo clique.} Para descobrir a torre clicada, o cursor é
levado ao NDC e desprojetado com $(PV)^{-1}$ nos planos near e far (com divisão
por $w$), o que gera um raio no mundo. O raio é testado contra as mesmas caixas
pelo método das faixas (\emph{slabs}), e a torre escolhida é a de menor
parâmetro $t$, isto é, a mais próxima da câmera. A reação é a pontuação: a torre
troca de lado e o corte é recontado.

% =====================================================================
\section{Capturas de tela}
% =====================================================================
As capturas foram geradas por \arq{relatorio/capturas.py}, que roda o jogo de
verdade, aciona as teclas pelos mesmos callbacks do jogador e lê o
\emph{backbuffer} com \texttt{ctx.screen.read()} antes de cada troca de buffers.
Como a leitura não inclui a decoração da janela, o texto da barra de título no
instante da captura foi escrito acima de cada imagem. A situação de vitória foi
preparada pelo script (a partição foi posta a uma inversão da meta) e concluída
pela regra do jogo (\texttt{inverter\_torre}). No canto superior esquerdo de
cada imagem está o placar: a barra amarela é a fração de arestas no corte, o
risco verde é a meta e a barra azul é o tempo restante.

\newcommand{\fig}[2]{\begin{subfigure}[t]{0.48\linewidth}\centering
  \includegraphics[width=\linewidth]{#1}\caption{#2}\end{subfigure}}

\begin{figure}[H]
\centering
\fig{inicio}{Início (\texttt{START}).}\hfill
\fig{jogando}{Jogo em andamento (\texttt{PLAYING}), câmera livre.}\\[6pt]
\fig{pausa}{Pausa (\texttt{PAUSED}).}\hfill
\fig{camera_aerea}{Câmera aérea.}
\caption{Estados do jogo e os dois modos de câmera.}
\end{figure}

\begin{figure}[H]
\centering
\fig{derrota}{Fim de jogo: derrota (tempo esgotado).}\hfill
\fig{vitoria}{Fim de jogo: vitória (corte atingiu a meta).}
\caption{Os dois finais.}
\end{figure}

\begin{figure}[H]
\centering
\fig{luz_direcional}{Só a luz direcional (\tecla{1}).}\hfill
\fig{luz_pontual}{Só a luz pontual móvel (\tecla{2}).}\\[6pt]
\fig{luz_choque_torres}{Só os choques das torres: fontes pontuais locais (\tecla{4}).}\hfill
\fig{luz_choque_arestas}{Só os choques das arestas: fontes lineares (\tecla{5}).}\\[6pt]
\fig{texturas}{Todas as luzes. Texturas: torre (metal, mapeamento cilíndrico), piso e bordas (repetidas).}
\caption{Cada tipo de fonte de luz e as texturas.}
\end{figure}

% =====================================================================
\section{Divisão de tarefas}
% =====================================================================

\begin{itemize}
\item \textbf{Lucas Ataide}: escolha do tema e da mecânica; \arq{grafo.py}
  (geração do grafo, corte aleatorizado e busca local); \arq{jogo.py} (estado da
  partida, cronômetro, sorteio de posições, seleção pelo clique e colisão);
  adaptação de \arq{camera.py} (modos e raio do cursor); \arq{geometria.py};
  conversão da torre para OFF; em \arq{main.py}, a janela, os callbacks, o laço,
  os estados, o desenho das arestas e dos anéis e o placar; script de capturas e
  relatório.
\item \textbf{Gabriel Frigo}: \arq{renderizador.py} (passadas de desenho, matriz
  normal, luzes no espaço da câmera, controle do sRGB); \arq{material.py} (luzes e
  materiais); integração de \arq{malha.py} e \arq{texturas.py}; shaders
  \arq{cena.vert/.frag} e \arq{linha.vert/.frag}; luzes dos choques (fontes locais
  nas torres e lineares nas arestas, teclas 3 a 5) e calibração da iluminação e dos
  materiais; escolha e integração das texturas do Poly Haven, incluindo a das
  bordas; reorganização de \arq{main.py} em torno do renderizador, resolução de
  conflito de merge e limpeza de código/comentários; lista de controles;
  \arq{Makefile}, \arq{requirements.txt} e \arq{README.md}.
\end{itemize}

% =====================================================================
\section{Declaração de uso de IA}
% =====================================================================
\begin{table}[H]
\centering\small
\begin{tabularx}{\linewidth}{@{}llY@{}}
\toprule
\textbf{Quem} & \textbf{Ferramenta} & \textbf{Uso} \\
\midrule
Lucas & Gemini (Google), em linha de comando &
Sugestão do tema; explicações de conceitos (DFS aleatorizada, iluminação,
envelope do shader); revisão do \arq{grafo.py} e do \arq{jogo.py}; código de
\arq{geometria.py} (piso, cubo, anel, segmento, leitura do OFF), dos modos da
câmera, de uma versão do \arq{jogo.py} e do desenho das arestas e anéis; conversão OBJ$\to$OFF;
\arq{.gitignore} e configuração do verificador de tipos; dois guias em PDF
(iluminação e envelope). \\
Lucas & Claude (Anthropic), Claude Code &
Explicações e revisão de código; código da seleção pelo clique (raio do cursor
e teste raio $\times$ caixa), do placar, da colisão da câmera e dos títulos de
fim de jogo; resolução de um conflito no git; script de capturas; redação deste
relatório a partir do código. \\
Gabriel & Gemini (Google) &
Análise de código morto, bugs e verificação de requisitos do enunciado;
implementação da camada gráfica a partir das especificações do projeto:
\arq{renderizador.py} (fachada OpenGL 4.0 Core, cache de VAOs, matriz normal,
luzes no espaço da câmera e controle de sRGB), shaders \arq{cena.vert/.frag} e
integração de \arq{linha.vert/.frag}, adaptação de \arq{material.py}, integração
de \arq{malha.py} em \arq{geometria.py} e reorganização de \arq{main.py} e
\arq{Makefile}; implementação das luzes dos choques (fontes pontuais locais nas
torres e lineares nas arestas, com atenuação de alcance finito e teclas
\tecla{3} a \tecla{5}); integração das texturas do Poly Haven (incluindo
\texttt{u\_uv\_escala} nas bordas), diagnóstico e calibração das intensidades de
luz e dos materiais; padronização de comentários para docstrings; resolução de
conflito de merge em \arq{main.py}; confirmação das licenças das texturas e
preenchimento desta declaração no relatório. \\
\bottomrule
\end{tabularx}
\end{table}

Os autores revisaram o material gerado com ajuda de IA e assumem total
responsabilidade técnica e intelectual por ele.

\end{document}
